\documentclass[11pt]{article}

\usepackage[a4paper,margin=1in]{geometry}
\usepackage{amsmath,amssymb,amsthm}
\usepackage{mathtools}
\usepackage{xcolor}
\usepackage{hyperref}
\usepackage{cleveref}
\usepackage{microtype}
\usepackage{enumitem}
\usepackage{booktabs}
\usepackage{tikz-cd}
\usepackage{tikz}
\usepackage{everypage}
\usepackage{listings}

\hypersetup{
  colorlinks=true,
  linkcolor=blue!50!black,
  citecolor=blue!50!black,
  urlcolor=blue!50!black
}

\theoremstyle{plain}
\newtheorem{theorem}{Theorem}[section]
\newtheorem{proposition}[theorem]{Proposition}
\newtheorem{lemma}[theorem]{Lemma}
\newtheorem{corollary}[theorem]{Corollary}

\theoremstyle{definition}
\newtheorem{definition}[theorem]{Definition}
\newtheorem{example}[theorem]{Example}

\theoremstyle{remark}
\newtheorem{remark}[theorem]{Remark}
\newtheorem{audit}[theorem]{Audit}

\lstdefinestyle{leanstyle}{
  basicstyle=\ttfamily\footnotesize,
  breaklines=true,
  columns=fullflexible,
  keepspaces=true,
  xleftmargin=1em,
  showstringspaces=false,
  frame=single,
  framesep=4pt,
  framerule=0.4pt,
  inputencoding=utf8,
  extendedchars=true,
  literate={Ω}{{$\Omega$}}1
           {≠}{{$\neq$}}1
           {∀}{{$\forall$}}1
           {∃}{{$\exists$}}1
           {→}{{$\to$}}1
           {↦}{{$\mapsto$}}1
           {↪}{{$\hookrightarrow$}}1
           {ℂ}{{$\mathbb{C}$}}1
           {ι}{{$\iota$}}1
           {α}{{$\alpha$}}1
           {ρ}{{$\rho$}}1
}
\lstset{style=leanstyle}

\newcommand{\HtH}{H^{2}}
\newcommand{\KB}{K_{B}}
\newcommand{\C}{\mathbb{C}}
\newcommand{\D}{\mathbb{D}}
\newcommand{\Z}{\mathbb{Z}}
\newcommand{\N}{\mathbb{N}}
\newcommand{\R}{\mathbb{R}}
\newcommand{\Q}{\mathbb{Q}}
\newcommand{\DQ}{\mathcal{D}_{Q}}
\newcommand{\Yon}{\mathsf{y}}
\newcommand{\code}[1]{\texttt{#1}}
\newcommand{\NoPhantom}{\mathrm{NoPhantom}}
\newcommand{\YBDC}{\mathrm{YBDC}}
\newcommand{\BBRDS}{\mathrm{BBRDS}}
\newcommand{\OFK}{\mathrm{OFK}}
\newcommand{\RH}{\mathrm{RH}_{\mathrm{classical}}}
\newcommand{\Obs}{\Omega}

\title{Classical RH in the No-Phantom Language: \\
       \large the Vol6 Synthesis and its Single Residual Obstruction \\
       \large (Volume VI, Paper 07)}
\author{Yoneda AI Research \\
        \texttt{research@yonedaai.com}}
\date{2026-05-07}

\begin{document}
\maketitle

\begin{abstract}
This is the synthesis paper of Volume~VI of the HoTT/Yoneda Riemann
hypothesis programme. Volume~V proved the conditional theorem
\(\code{RH\_classical\_of\_no\_phantom\_language\_breakthrough} \;:\;
\code{NoPhantomLanguageBreakthrough} \to \RH\), reducing classical
RH to two named payloads: an off-critical defect kernel
(Target~1) and a Yoneda/Blaschke detector calculus (Target~2).
Volume~VI's proof sprint produced six papers, each shipping a Lean
module, four of which terminate at named obstructions because Vol5
ships an \emph{opaque} \(\code{burnolBlaschkeFactor.modelSpaceCarrier}\)
with no introduction rule plus three opaque admissibility atoms
\(\code{BlaschkeDefectDualizable}\), \(\code{BlaschkeDefectPolarized}\),
and \(\code{BlaschkeDefectRegularized}\).

The honest synthesis state of Vol6 is therefore
\emph{not} an unconditional proof of classical RH. We instead deliver:

\begin{itemize}[leftmargin=2em]
  \item a wrapper around the Vol5 master theorem under the canonical
        synthesis name
        (\code{RH\_classical\_new\_language\_of\_payload});
  \item a substantive parameterised theorem
        \code{RH\_classical\_new\_language\_of\_inputs} which builds the
        breakthrough payload from a Paper~03 transport bridge and a
        Paper~06 sub-target bundle;
  \item the single canonical residual obstruction
        \code{Vol6FinalObstruction} that collects exactly the four named
        introduction-rule fields produced by Papers 03--06; and
  \item the principal theorem
        \code{RH\_classical\_new\_language\_of\_obstruction}, which
        derives \(\RH\) from any inhabitant of
        \code{Vol6FinalObstruction} together with an explicit
        non-opaque detector and \code{DQObj} witness.
\end{itemize}

We deliberately do \emph{not} ship an unconditional
\code{RH\_classical\_new\_language : RH\_classical}: producing such a
term inside Vol6 without a forbidden
\code{sorry}/\code{axiom}/\code{opaque} declaration is impossible
given the structural obstructions of the Vol5 surface.  We perform a
verbatim audit against the four success-criteria checks of
\code{next-steps.md}, run the Vol6 build, count
\code{sorry}/\code{admit}/\code{axiom}/\code{opaque} occurrences across
\code{lean/Vol6/}, and read off the
\code{\#print axioms} dependency list of the synthesis principal
theorem. All four success criteria pass except the
unconditional-closure check, which we explicitly re-state as the
Vol6 \(\Obs\) obstruction.

The output is a clean, honest reduction of classical RH to a single
\(\Type 1\)-valued obstruction structure, ready for Vol7.
\end{abstract}

\tableofcontents
\newpage

\section{Introduction}
\label{sec:introduction}

\subsection{The Vol5 conditional reduction, restated}
\label{sec:vol5-restate}

The starting point of Vol6 is the Vol5 master theorem
\begin{equation}
  \label{eq:rh-of-breakthrough}
  \code{RH\_classical\_of\_no\_phantom\_language\_breakthrough}
    \;:\;
    \code{NoPhantomLanguageBreakthrough} \;\longrightarrow\;
    \RH,
\end{equation}
proved unconditionally inside Vol5
(\code{Vol5.YonedaNymanTrackA.NoPhantomLanguage}, line~137). Its
hypothesis is the record
\begin{equation}
  \label{eq:breakthrough-record}
  \code{NoPhantomLanguageBreakthrough} \;=\;
    \begin{pmatrix}
      \code{detectorCalculus} & : & \code{YonedaBlaschkeDetectorCalculus} \\
      \code{offCriticalKernel} & : & \code{OffCriticalZeroDefectKernel}
    \end{pmatrix},
\end{equation}
where \(\code{YonedaBlaschkeDetectorCalculus}\) wraps the canonical RKHS
detector semantics
\(\BBRDS = \code{BurnolBlaschkeRKHSDetectorSemantics}\), itself a
five-component record
\begin{equation}
  \label{eq:bbrds}
  \BBRDS \;=\;
    (\,\mathrm{detector},\,\mathrm{rational},\,\mathrm{externalization},\,
      \mathrm{orthogonality},\,\mathrm{admissibility}\,),
\end{equation}
and \(\code{OffCriticalZeroDefectKernel}\) is the two-line proposition
\begin{equation}
  \label{eq:offCriticalKernel}
  \code{OffCriticalZeroDefectKernel} \;=\;
    \bigl(\,\code{ExistsOffCriticalZero}
      \to \code{Nonempty}\,\code{burnolBlaschkeFactor.modelSpaceCarrier}\,\bigr).
\end{equation}

The mathematical pipeline \eqref{eq:rh-of-breakthrough} records that
classical RH follows from these two payloads via the Vol5 chain
\[
  \begin{array}{l}
    \text{off-critical zero}
      \xrightarrow{\text{kernel}} \text{nonzero }
        \code{burnolBlaschkeFactor.modelSpaceCarrier} \\
    \xrightarrow{\text{detector}} \text{nonzero defect detected
      by RKHS probe}\\
    \xrightarrow{\text{externalize}} \text{Yoneda representable in
      }\DQ\\
    \xrightarrow{\text{orthogonality}} \text{contradicts cokernel
      invisibility}\\
    \xrightarrow{\text{admissibility}}\Longrightarrow \text{no
      off-critical zero exists}
    \;\Longrightarrow\;\RH.
  \end{array}
\]

The Vol6 sprint's job is to inhabit \eqref{eq:breakthrough-record}.

\subsection{The honest Vol6 sprint inventory}
\label{sec:honest-inventory}

Six prior Vol6 papers (and their Lean modules) carry out this
inhabitation as far as the Vol5 opaque/axiom surface permits. The
\emph{honest} synthesis state is summarised in
\Cref{tab:honest-inventory}.

\begin{table}[h]
\centering
\small
\begin{tabular}{@{}l l l l@{}}
\toprule
Paper & Lean module
      & Closed unconditionally?
      & Obstruction module \\
\midrule
01 & \code{Vol6.FiniteBlaschkePacket}
   & YES
   & none \\
02 & \code{Vol6.FiniteRankDetector}
   & YES
   & \code{Obstruction.FiniteRankDetectorObstruction} (canonical lift) \\
03 & \code{Vol6.OffCriticalDefectKernel}
   & NO --- parameterised by
     \code{OffCriticalDefectKernelBridge}
   & \code{Obstruction.OffCriticalDefectKernelObstruction} (RF-01) \\
04 & \code{Vol6.RationalDilationExternalization}
   & NO --- parameterised by
     \code{ExternalizationIntroductionRule}
   & \code{Obstruction.RationalDilationExternalizationObstruction} (RF-02) \\
05 & \code{Vol6.QuotientOrthogonalityAdmissibility}
   & NO --- parameterised by two introduction packages
   & \code{Obstruction.QuotientOrthogonalityObstruction} +
     \code{Obstruction.AdmissibilityObstruction} (RF-02, RF-03) \\
06 & \code{Vol6.YonedaBlaschkeDetectorCalculus}
   & NO --- parameterised by \code{BundledSubTargets}
   & \code{Obstruction.YonedaBlaschkeDetectorCalculusObstruction} \\
\bottomrule
\end{tabular}
\caption{The honest Vol6 sprint inventory. Two papers (01 and 02)
close their principal theorems unconditionally; four papers terminate
at named introduction-rule obstructions. The structural reason is
documented in \Cref{sec:rf01-rf04}.}
\label{tab:honest-inventory}
\end{table}

The structural reason that Papers~03--06 cannot close unconditionally
is \emph{not} a missing finite-dimensional argument. It is that the
Vol5 surface ships
\begin{itemize}[leftmargin=2em]
  \item \code{axiom burnolBlaschkeFactor : BurnolBlaschkeFactor}
        (\code{Vol5.YonedaNymanTrackA.BurnolCore}, line~42), whose
        \(\code{modelSpaceCarrier} : \Type\) field has no constructor,
        injection, or identification with any concrete
        \(\Type\); and
  \item \code{opaque BlaschkeDefectDualizable}, \code{Polarized},
        \code{Regularized} (\code{Vol5.YonedaNymanTrackA.NoPhantomBlaschkeDefect},
        lines~66--70), and \code{opaque
        DefectDetectedByRationalDilationRepresentable}
        (line~28), each declared as \(\code{opaque} \cdot \to \cdot
        \to \Prop\) with no introduction rule.
\end{itemize}
Both facts are \emph{deliberate} risk flags of the Vol6 brief, recorded
as RF-01 through RF-04 of \code{.knowledge-base.md}~\S8.

\subsection{Synthesis goal: the single residual obstruction}
\label{sec:synthesis-goal}

Given \Cref{tab:honest-inventory}, the Vol6 synthesis can take one of
two attitudes:

\begin{enumerate}
  \item \emph{Pretend}: introduce a new \(\code{axiom}\) or
        \(\code{opaque}\) closing the obstruction, and ship a
        formally-typeable but mathematically dishonest term
        \(\code{RH\_classical\_new\_language : RH\_classical}\). This
        is forbidden by the Vol6 hard rules.
  \item \emph{Be honest}: assemble the four named introduction-rule
        fields produced by Papers 03--06 into a \emph{single} canonical
        residual obstruction structure
        \(\code{Vol6FinalObstruction}\), prove that \(\RH\) follows
        from any inhabitant of \(\code{Vol6FinalObstruction}\)
        together with non-opaque detector and \code{DQObj}
        data, and document
        \(\code{Vol6FinalObstruction}\) as the canonical Vol7
        deliverable.
\end{enumerate}

The present paper takes attitude (2). The Lean module
\code{Vol6.RHClassicalNewLanguage} ships exactly the parameterised
synthesis theorems --- \emph{nothing more} --- and the
\(\code{\#print\;axioms}\) audit confirms that the principal theorem
inherits only Mathlib foundational axioms plus the pre-existing Vol5
axiom surface; no RH, no Nyman density, no Beurling-Nyman, no
RH-equivalent appears.

\subsection{Outline of the paper}

\Cref{sec:inventory} reviews each of the six prior Vol6 papers and
states precisely what each closed and what each parameterised.
\Cref{sec:rf01-rf04} states the four risk flags
RF-01--RF-04 explicitly. \Cref{sec:composition} defines the assembled
structure \(\code{Vol6FinalObstruction}\) and proves the
principal synthesis theorem. \Cref{sec:audit} performs the
verbatim success-criteria audit from \code{next-steps.md}, including
the lake build, the no-axiom grep, and the
\(\code{\#print\;axioms}\) check. \Cref{sec:verdict} states the
honest synthesis verdict. \Cref{sec:vol7-roadmap} sketches the four
bridge constructions Vol7 must produce. Appendices reproduce the
key Lean signatures.

\subsection{Notation and conventions}

We use the following conventions throughout. The complex unit disc is
\(\D = \{z \in \C : |z| < 1\}\); its boundary is the unit circle
\(\partial \D = \{z \in \C : |z| = 1\}\). The Hardy space \(\HtH\) is the
classical \(L^{2}(\partial \D)\)-completion of bounded analytic functions
on \(\D\) under the boundary-trace inner product
\[
  \langle f, g \rangle_{\HtH}
    \;=\; \int_{0}^{2\pi}
      f(e^{i\theta})\,\overline{g(e^{i\theta})}\,
      \frac{d\theta}{2\pi}.
\]
A Blaschke product over a finite zero set
\(\{\alpha_{i}\}_{i \in I} \subset \D\) is
\[
  B_{P}(z) \;=\; \prod_{i \in I}
    \frac{z - \alpha_{i}}{1 - \overline{\alpha_{i}}\,z},
\]
and its associated \emph{model space} is
\[
  K_{B_{P}} \;=\; \HtH \,\ominus\, B_{P}\,\HtH
   \;=\; \HtH \,/\, B_{P}\,\HtH,
\]
which is \(\C\)-linear of dimension \(|I|\) and has Pick / Cauchy basis
\(\{(1 - \overline{\alpha_{i}}\,z)^{-1}\}_{i \in I}\).
On the Lean side we work with the surrogate
\(P.\code{ZeroIndex} \to \C\), which is the coefficient space of this
basis under the standard isomorphism.

\paragraph{Vol5 notation.}
Throughout we denote by \(\code{burnolBlaschkeFactor}\) the
canonical \(\code{BurnolBlaschkeFactor}\) supplied as an axiom in
\code{Vol5.YonedaNymanTrackA.BurnolCore}, and by
\(\code{blaschkeDefectObject}\) its packaging as a
\(\code{BlaschkeDefectObject}\) in
\code{Vol5.YonedaNymanTrackA.BlaschkeDefectObject}.  The associated
condensed Hilbert defect is
\(\code{burnolBlaschkeCondensedHilbertDefect}\) defined in
\code{Vol5.YonedaNymanTrackA.CondensedHilbertDefect}.  The opaque
Vol5 atoms governing detection and admissibility live in
\code{Vol5.YonedaNymanTrackA.NoPhantomBlaschkeDefect}.  The
master conditional theorem we apply is
\code{RH\_classical\_of\_no\_phantom\_language\_breakthrough} of
\code{Vol5.YonedaNymanTrackA.NoPhantomLanguage}, line~137.

\paragraph{Lean sigil conventions.}
Throughout the paper we treat code blocks as \emph{verbatim} copies of
the corresponding Lean term, modulo whitespace.  Where space requires
abbreviation we write the abbreviation in full at first appearance
(e.g.\ ``\code{burnolBlaschkeCondensedHilbertDefect}'' is sometimes
written in shorthand inside record literals as \code{K}).  We use
\(\Type 1\) to denote Lean's first-universe-above-Prop, in which all
data-carrying structures of this paper live.

\section{Inventory of the six Vol6 papers}
\label{sec:inventory}

We summarise each prior Vol6 paper from the synthesis viewpoint:
which Lean theorem it ships, what is closed, and what is
parameterised.

\subsection{Paper 01 --- FiniteBlaschkePacket}
\label{sec:p01}

\paragraph{Module.}
\code{Vol6.FiniteBlaschkePacket}.

\paragraph{Status.}
Closed unconditionally. The principal theorem
\begin{lstlisting}
theorem finite_packet_model_space_nonempty
    (P : FiniteBlaschkePacket) (h : Nonempty P.ZeroIndex) :
    Nonempty (finiteBlaschkeModelSpace P).carrier
\end{lstlisting}
is proved by direct construction. The carrier of
\(\code{finiteBlaschkeModelSpace}\) is \(P.\code{ZeroIndex} \to \C\)
and the witness is the constant-zero function. The hypothesis
\(\code{Nonempty}\,P.\code{ZeroIndex}\) is recorded for spec
contract reasons. Auxiliary lemmas
\code{finitePacketPickVector\_self} and \code{finitePacketPickVector\_other}
expose the Pick basis used by Paper~02.

\paragraph{What it provides to synthesis.}
The finite-dimensional model-space surrogate, with a constructive
non-emptiness witness.  No obstruction module is paired with Paper~01:
nothing about it depends on the Vol5 opaque surface.

\subsection{Paper 02 --- FiniteRankDetector}
\label{sec:p02}

\paragraph{Module.}
\code{Vol6.FiniteRankDetector}.

\paragraph{Status.}
Closed unconditionally at the finite-packet level; canonical lift
gated. The principal theorem is
\begin{lstlisting}
theorem finite_rank_detector_complete
    (P : FiniteBlaschkePacket) (h : Nonempty P.ZeroIndex) :
    forall v : (finiteBlaschkeModelSpace P).carrier,
      v ≠ (fun _ => (0 : Complex)) ->
        ∃ i : P.ZeroIndex, evalKernel P i v
\end{lstlisting}
proved by function-extensionality on the coordinate space, which
captures the Cauchy/Pick determinant nondegeneracy at the finite-rank
abstraction layer.

\paragraph{Obstruction module.}
\code{Vol6.Obstruction.FiniteRankDetectorObstruction} records the
\emph{canonical-detector lift} obstruction:
\code{FiniteRankDetectorCompletionDatum} carries a transport
\(\bigl((\code{finiteBlaschkeModelSpace}\,P).\code{carrier}\bigr) \to
   \code{burnolBlaschkeFactor.modelSpaceCarrier}\) plus a coherence
field. No inhabitant is supplied, in compliance with RF-01.

\paragraph{What it provides to synthesis.}
The detector data structure used by Paper~04 and Paper~06 is built
from this paper's combinatorics. The canonical-lift obstruction is
absorbed by Paper~03's bridge \(\code{OffCriticalDefectKernelBridge}\)
in the synthesis structure (Section~\ref{sec:composition}).

\subsection{Cross-paper composition lemma (P01 + P02)}

The two unconditionally-closed papers compose as follows.  P01 supplies
\(\code{(finiteBlaschkeModelSpace}\,P).\code{carrier}
   = P.\code{ZeroIndex} \to \C\) and the constant-zero witness; P02
supplies the coordinate-detection predicate
\(\code{evalKernel}\,P\,i\,v = (v\,i \neq 0)\) and the
function-extensionality lemma
\code{coordinate\_witness\_of\_nonzero}.  Together they give the
finite-rank version of Target~1 that Paper~03 transports to the
canonical level:

\begin{lemma}[Finite-rank visibility, paraphrased]
\label{lem:finite-rank-visibility}
For every finite Blaschke packet \(P\) and every nonzero vector
\(v \in (\code{finiteBlaschkeModelSpace}\,P).\code{carrier}\),
some Pick-basis evaluation functional \(\code{evalKernel}\,P\,i\)
detects \(v\), where ``detect'' means
\(v_{i} \neq 0\).
\end{lemma}

The proof is simply
\code{finite\_rank\_detector\_complete}.  Mathematically, this is the
finite-dimensional Cauchy/Pick determinant nondegeneracy
\[
  \det\,\bigl(\,(K(\alpha_{i}, \alpha_{j}))_{i,j}\,\bigr)
   \;=\;
   \frac{\prod_{i < j} |\alpha_{i} - \alpha_{j}|^{2}}
        {\prod_{i, j}
           (1 - \alpha_{j} \overline{\alpha_{i}})}
   \neq 0
\]
which makes the Pick-basis a separating family for \(K_{B_{P}}\)
when the \(\alpha_{i}\) are distinct.  In Lean we work with the
weaker statement that the standard basis of \(P.\code{ZeroIndex}\to \C\)
is separating, which is just function-extensionality.

\subsection{Paper 03 --- OffCriticalDefectKernel}
\label{sec:p03}

\paragraph{Module.}
\code{Vol6.OffCriticalDefectKernel}.

\paragraph{Status.}
\emph{Conditional.} The principal theorem shipped is
\begin{lstlisting}
theorem off_critical_zero_defect_kernel_of_bridge
    (bridge : OffCriticalDefectKernelBridge) :
    OffCriticalZeroDefectKernel
\end{lstlisting}
proved by chaining the Möbius parameter map
\(\alpha(s) = (s-1)/(s+1)\), the singleton finite Blaschke packet
\(P_{s} = \{\alpha(s)\}\), Paper~01's
\(\code{finite\_packet\_model\_space\_nonempty}\), and the bridge's
transport map.  No unconditional
\code{off\_critical\_zero\_defect\_kernel} is shipped.

\paragraph{Obstruction module.}
\code{Vol6.Obstruction.OffCriticalDefectKernelObstruction} states
\(\code{OffCriticalDefectKernelBridge}\), a \(\Type 1\)-valued
structure carrying a single field
\(\code{transport} : \forall P,\; \code{Nonempty}\,P.\code{ZeroIndex}\,\to\,
   (\code{finiteBlaschkeModelSpace}\,P).\code{carrier} \to
   \code{burnolBlaschkeFactor.modelSpaceCarrier}\),
i.e.\ exactly the missing Vol5 transport (RF-01).

\subsection{Paper 04 --- RationalDilationExternalization}
\label{sec:p04}

\paragraph{Module.}
\code{Vol6.RationalDilationExternalization}.

\paragraph{Status.}
\emph{Conditional.} The principal definition
\begin{lstlisting}
noncomputable def burnol_blaschke_detector_externalization
    (S : RKHSModelSpaceDetector burnolBlaschkeCondensedHilbertDefect)
    (X : DQObj)
    (iota :
      ExternalizationIntroductionRule
        burnolBlaschkeCondensedHilbertDefect) :
    RKHSDetectorExternalizationToRationalRepresentables
      burnolBlaschkeCondensedHilbertDefect S
      (canonicalRationalDilationSemantics X)
\end{lstlisting}
ships the externalization given an introduction rule \(\iota\) for the
opaque \(\code{DefectDetectedByRationalDilationRepresentable}\)
predicate.  An auxiliary unconditional inhabitant
\code{burnol\_blaschke\_trivial\_externalization} handles the
\(\code{InternalBlaschkeTriviality}\) case using only the empty detector.

\paragraph{Obstruction module.}
\code{Vol6.Obstruction.RationalDilationExternalizationObstruction}
states \(\code{ExternalizationIntroductionRule}\;K\), a one-field
\(\Type 1\)-valued structure
\(\code{intro} : \forall F,\;\code{RationalDilationRepresentable}\,F
   \to \code{CondensedHilbertDefectDetectedByRationalDilation}\,K\,F\).
This is RF-02.

\subsection{Paper 05 --- QuotientOrthogonalityAdmissibility}
\label{sec:p05}

\paragraph{Module.}
\code{Vol6.QuotientOrthogonalityAdmissibility}.

\paragraph{Status.}
\emph{Conditional, two introductions.} The two principal
deliverables are
\begin{lstlisting}
noncomputable def burnol_blaschke_quotient_orthogonality_of_intro
    (Q : CanonicalQuotientOrthogonalityIntroduction) :
    QuotientOrthogonalityInvisibility blaschkeDefectObject

theorem burnol_blaschke_defect_is_admissible_of_intro
    (A : CanonicalAdmissibilityIntroduction) :
    BurnolBlaschkeDefectIsAdmissible
\end{lstlisting}
The cokernel construction takes
\(\code{OrthogonalToRepresentable}\,F :=
   \code{YonedaNymanRepresentableImage}\,F\) (the canonical
cokernel-orthogonality choice). Both
deliverables are also re-stated in their iff-form
\code{burnol\_blaschke\_quotient\_orthogonality\_iff} and
\code{burnol\_blaschke\_defect\_is\_admissible\_iff}, witnessing
that the introduction packages are precisely as
strong as their targets.

\paragraph{Obstruction modules.}
\code{Vol6.Obstruction.QuotientOrthogonalityObstruction} (RF-02
partial; one cokernel-introduction field) and
\code{Vol6.Obstruction.AdmissibilityObstruction} (RF-03; three
fields, one per opaque admissibility atom).

\subsection{Mathematical content of Paper 04 (the externalization picture)}

Mathematically, Paper~04 is the statement that any RKHS detector for
the canonical Burnol/Blaschke defect is represented by a Yoneda
\(\DQ\)-presheaf.  The conceptual picture is
\[
  \begin{tikzcd}[column sep=large]
    \text{nonzero }v \in K_{B}
      \arrow[r, "{\,\text{detector}\,}"]
      \arrow[rr, bend right=12, "{\,\text{externalization}\,}"', dashed]
    & d \in S.\code{Detector}
      \arrow[r, "{\,\Yon \circ R.\code{objectOf}\,}"]
    & F \in \DQ\text{-presheaf}
  \end{tikzcd}
\]
where the dashed arrow is the externalization that Paper~04 is
asked to supply.  Concretely, the Vol6 implementation takes
\(R = \code{canonicalRationalDilationSemantics}\,X\) for a fixed
\code{DQObj} \(X\) (the parameter type is \code{Unit}, so the object
is constant) and \(F = \Yon\,X\) for every detector index.  The
\code{presheaf\_eq} field is then reflexive; the \code{detected\_of\_detector}
field uses the introduction rule \(\iota\) to transport from
\(\Yon\,X\) being a \code{RationalDilationRepresentable} to the
condensed Hilbert defect being detected.

\subsection{Paper 06 --- YonedaBlaschkeDetectorCalculus}
\label{sec:p06}

\paragraph{Module.}
\code{Vol6.YonedaBlaschkeDetectorCalculus}.

\paragraph{Status.}
\emph{Assembly only.} The principal definition
\begin{lstlisting}
def yoneda_blaschke_detector_calculus_of_subtargets
    (B : BundledSubTargets) :
    YonedaBlaschkeDetectorCalculus
\end{lstlisting}
is proved by a single record literal in the four-component bundle
type \code{BundledSubTargets} (detector, rational, externalization,
orthogonality, admissibility). The companion
\begin{lstlisting}
theorem RH_classical_of_subtargets_and_off_critical_kernel
    (B : BundledSubTargets) (K : OffCriticalZeroDefectKernel) :
    RH_classical
\end{lstlisting}
is the conditional master theorem in Paper~06's form.  No
unconditional Target~2 theorem is shipped.

\paragraph{Obstruction module.}
\code{Vol6.Obstruction.YonedaBlaschkeDetectorCalculusObstruction}
records the four upstream barriers in a string-data
\code{DetectorCalculusBarrier} structure (one
\code{UpstreamBarrier} per missing principal theorem).

\subsection{Mathematical content of Paper 05}

Mathematically, Paper~05 makes two distinct contributions, one
qualitative (orthogonality) and one quantitative (admissibility).

\paragraph{Orthogonality (Section A).}
The cokernel construction \(K_{B} = \HtH / B\,\HtH\) immediately
implies that any element of \(B\,\HtH\) is annihilated by the
canonical projection \(\pi : \HtH \to K_{B}\).  In particular,
the realisations of the Yoneda/Nyman representable presheaves --- which
factor through \(B\,\HtH\) --- are invisible to detection by
\(K_{B}\).  This is the only mathematics needed for orthogonality;
no Nyman density is invoked.  The Lean cokernel-introduction package
takes
\[
  \code{OrthogonalToRepresentable}\,F
   \;:=\;
   \code{YonedaNymanRepresentableImage}\,F,
\]
which is the canonical choice making \code{orthogonal\_of\_yoneda\_image}
hold reflexively and \code{no\_detection\_of\_orthogonal} reduce to
the cokernel rule.

\paragraph{Admissibility (Section B).}
At the finite-rank level, all three admissibility atoms reduce to
standard finite-dimensional linear algebra:
\begin{itemize}[leftmargin=2em]
  \item \emph{Dualizability:} the model space \(K_{B_{P}}\)
        \(=\) \(P.\code{ZeroIndex} \to \C\) is a finite-dimensional
        Hilbert module with the obvious self-pairing
        \((u, v) \mapsto \sum_{i} u_{i} \overline{v_{i}}\); the
        evaluation/coevaluation maps satisfy the snake identities
        because \(K_{B_{P}}\) is reflexive.
  \item \emph{Polarisation:} the Hardy inner product restricted to
        \(K_{B_{P}}\) is a non-degenerate Hermitian form
        \(\langle u, v \rangle = u^{*} G v\) where \(G\) is the
        Cauchy/Pick matrix; the polarisation identity
        \(\|u + v\|^{2} - \|u - v\|^{2}
           + i\|u + iv\|^{2} - i\|u - iv\|^{2} = 4\,\Re\langle u, v\rangle\)
        is automatic in any complex inner product space.
  \item \emph{Regularised convergence:} the resolvent family
        \((R(z))_{z \in \rho(T_{B_{P}})}\) of the truncated
        multiplication operator on \(K_{B_{P}}\) is a finite-rank
        rational family in \(z\); regularised convergence under
        spectral truncation is automatic.
\end{itemize}
The \emph{lift} from finite-rank to the canonical condensed defect
is the obstruction; this is precisely RF-03.

\subsection{Synthesis takeaway}

The two unconditionally-closed papers (01, 02) carry the
\emph{finite-rank} content; the four conditionally-closed papers
(03, 04, 05, 06) each isolate the precise opaque/axiom Vol5 surface
they cannot bypass.  The synthesis is therefore: collect the four
named introduction-rule fields into one structure, prove
\(\RH\) from any inhabitant, and document the structure as the
single canonical Vol7 deliverable.

\subsection{The dependency graph of Vol6}

Vol6 has a strict (DAG) dependency structure between its modules,
which we list compactly:
\begin{itemize}[leftmargin=2em]
  \item P01 \(\to\) \{P02, P03, P04, P05\}.
  \item P02 \(\to\) \{P03, P04, P06\}.
  \item P03 \(\to\) P07.
  \item P04 \(\to\) P06.
  \item P05A, P05B \(\to\) P06.
  \item P06 \(\to\) P07.
  \item Each obstruction module
        \texttt{Vol6.Obstruction.X} is imported by its paired
        proof module \texttt{Vol6.X}.
  \item P07 (this paper) imports all six prior proof modules
        and all six obstruction modules.
\end{itemize}
Edges denote direct Lean \code{import} relationships in
\code{lean/Vol6/}.  The synthesis paper P07 imports all six prior
papers and all six obstruction modules; the obstruction modules form
a parallel DAG in \code{lean/Vol6/Obstruction/}.

\paragraph{Reading order of the present paper.}
A reader familiar with Vol5 may proceed directly to
\Cref{sec:composition}.  A reader new to the Vol5 surface
should first read \Cref{sec:rf01-rf04} which states the
opaque/axiom obligations explicitly.  A reader interested only in
the audit may read \Cref{sec:audit}.  A reader planning the Vol7
work should focus on \Cref{sec:vol7-roadmap}.

\section{The Vol5 risk flags RF-01--RF-04, restated}
\label{sec:rf01-rf04}

For self-containedness we restate the four risk flags driving
\Cref{tab:honest-inventory}.

\paragraph{RF-01.}
\code{burnolBlaschkeFactor : BurnolBlaschkeFactor} is an
\code{axiom} in
\code{Vol5.YonedaNymanTrackA.BurnolCore} (line~42), with
\code{modelSpaceCarrier : Type} as a structure field. There is
no Vol5-exposed equality, injection, or constructor whose
conclusion lands in \code{burnolBlaschkeFactor.modelSpaceCarrier}.
Paper~03's bridge fills exactly this gap.

\paragraph{RF-02.}
\code{DefectDetectedByRationalDilationRepresentable :
   BlaschkeDefectObject \to DQPresheaf \to Prop} is
\code{opaque} in \code{Vol5.YonedaNymanTrackA.NoPhantomBlaschkeDefect}
(line~28), with no Vol5-exposed introduction rule. Paper~04's
\code{ExternalizationIntroductionRule} and Paper~05A's
cokernel-introduction package each fill one face of this gap.

\paragraph{RF-03.}
\code{BlaschkeDefectDualizable / Polarized / Regularized :
   BlaschkeDefectObject \to Prop} are all \code{opaque}
(lines~66--70), with no introduction rule. Paper~05B's three-field
\code{AdmissibilityIntroduction} package collects them.

\paragraph{RF-04.}
\code{HardySpaceH2 : Type} and \code{HardySubmodule : Type}
are \code{axiom} types in \code{BurnolCore} (lines~19--21), so
no Mathlib-backed inner-product computation can be transported into
the canonical model-space without the Paper~03 bridge of RF-01.
Paper~01 sidesteps this by working in the explicit type
\(P.\code{ZeroIndex} \to \C\); the canonical lift remains gated.

\subsection{Why the structural obstructions are not accidental}

It is worth pausing on \emph{why} Vol5 ships its key types and atoms
opaquely.  The Vol5 architecture deliberately separates the
\emph{interface} of the no-phantom language (which is what Vol5
proves: a contradiction emerges between admissibility and
phantom-defect rigidity) from the \emph{implementation} of the
canonical Burnol/Blaschke factor (which is a delicate piece of
analysis touching on \(L^{2}\) Mellin transforms, fractional-part
kernels, and the Beurling rigidity criterion).  The opaque/axiom
declaration is a deliberate \emph{interface freeze}: any sound
implementation that satisfies the opaque declarations' types
discharges the Vol5 chain.  Vol6 is the first sprint to attempt
that implementation; its honest verdict is that the four
introduction rules listed above are exactly the missing
implementation work.

\subsection{The obstructions are downward-irreducible}

Each of the four fields of \(\code{Vol6FinalObstruction}\)
(\Cref{def:final-obstruction}) is downward-irreducible in the
following sense: removing it from \(\Obs\) and trying to construct
\(\RH\) from the remaining three would require either reproducing
that field's content from the other three (impossible: the four
obstructions name \emph{disjoint} pieces of the Vol5 opaque
surface) or proving an unrelated stronger lemma that subsumes it.
The \emph{iff} audits in
\code{Vol6.Obstruction.QuotientOrthogonalityObstruction} and
\code{Vol6.Obstruction.AdmissibilityObstruction}
(theorems
\code{quotient\_orthogonality\_invisibility\_iff\_introduction} and
\code{admissibility\_intro\_iff} respectively) prove that each
introduction package is \emph{exactly} as strong as its target
structure --- not stronger, not weaker.  The same holds for the
externalization (Paper~04) and the canonical-lift bridge
(Paper~03).

\section{Composition: the assembled obstruction and its discharge}
\label{sec:composition}

We now state the central synthesis composition.  All Lean signatures
in this section are reproduced verbatim from
\code{lean/Vol6/RHClassicalNewLanguage.lean}.

\subsection{Wrapper around the Vol5 master theorem}

\begin{theorem}[Wrapper]
\label{thm:wrapper}
The function
\begin{lstlisting}
theorem RH_classical_new_language_of_payload
    (payload : NoPhantomLanguageBreakthrough) :
    RH_classical :=
  RH_classical_of_no_phantom_language_breakthrough payload
\end{lstlisting}
is a one-line wrapper around \eqref{eq:rh-of-breakthrough}. It is
provable in Vol6 because the Vol5 master theorem is unconditional.
\end{theorem}

\subsection{Substantive deliverable: RH from Paper~03 + Paper~06}

\begin{theorem}[RH from inputs]
\label{thm:rh-from-inputs}
Given a Paper~03 transport bridge and a Paper~06 sub-target bundle,
classical RH follows:
\begin{lstlisting}
theorem RH_classical_new_language_of_inputs
    (P03 : OffCriticalDefectKernelBridge)
    (P06 : Vol6.BundledSubTargets) :
    RH_classical :=
  RH_classical_of_no_phantom_language_breakthrough
    { detectorCalculus :=
        Vol6.assembleYonedaBlaschkeDetectorCalculus P06,
      offCriticalKernel :=
        Vol6.OffCriticalDefectKernelNS
          .off_critical_zero_defect_kernel_of_bridge P03 }
\end{lstlisting}
\end{theorem}

The proof body assembles the breakthrough payload field-wise from
Paper~03 (\code{off\_critical\_zero\_defect\_kernel\_of\_bridge}) and
Paper~06 (\code{assembleYonedaBlaschkeDetectorCalculus}) and then
applies \eqref{eq:rh-of-breakthrough}.

\subsection{The single canonical residual obstruction}

\begin{definition}[Vol6 final obstruction]
\label{def:final-obstruction}
\begin{lstlisting}
structure Vol6FinalObstruction : Type 1 where
  bridgeRF01    : OffCriticalDefectKernelBridge
  externalIntro : ExternalizationIntroductionRule
                    burnolBlaschkeCondensedHilbertDefect
  cokernelIntro : CanonicalQuotientOrthogonalityIntroduction
  admissIntro   : CanonicalAdmissibilityIntroduction
\end{lstlisting}
This is the assembled structure collecting exactly the four named
introduction-rule fields shipped by Papers~03--06.
\end{definition}

\begin{remark}
Field map to risk flags:
\begin{itemize}
  \item \code{bridgeRF01}: discharges RF-01 by transporting the
        finite-packet carrier into
        \code{burnolBlaschkeFactor.modelSpaceCarrier}.
  \item \code{externalIntro}: discharges RF-02 by introducing
        \code{CondensedHilbertDefectDetectedByRationalDilation}.
  \item \code{cokernelIntro}: the Paper-05A cokernel-introduction
        package supplying the Yoneda-image-not-detected
        clause.
  \item \code{admissIntro}: discharges RF-03 by introducing the
        three opaque admissibility atoms.
\end{itemize}
\end{remark}

\subsection{The principal synthesis theorem}

\begin{theorem}[Synthesis principal theorem]
\label{thm:synthesis-principal}
\begin{lstlisting}
theorem RH_classical_new_language_of_obstruction
    (Ω : Vol6FinalObstruction)
    (S : RKHSModelSpaceDetector burnolBlaschkeCondensedHilbertDefect)
    (X : DQObj) :
    RH_classical :=
  RH_classical_new_language_of_inputs Ω.bridgeRF01
    (bundledSubTargets_of_obstruction Ω S X)
\end{lstlisting}
The arguments \(S\) and \(X\) are explicit because they are
\emph{not} opaque-Prop discharge data: an
\(\code{RKHSModelSpaceDetector}\) inhabitant is suppliable by
the empty-detector under \(\code{InternalBlaschkeTriviality}\)
(\code{Vol6.Rational\-Dilation\-Externalization.burnol\-Blaschke\-Empty\-Detector}),
and a \code{DQObj} is suppliable by the Vol5 axiom-typed
\code{NymanKernel}.  Neither argument adds an obstruction beyond
\(\Obs\).
\end{theorem}

\begin{proof}[Lean proof]
The proof is a single-line composition.
\code{bundledSubTargets\_of\_obstruction} is constructed as a record
literal whose five fields are extracted from \(\Obs\):

\begin{itemize}
  \item \(\code{detector} := S\);
  \item \(\code{rational} :=
            \code{canonicalRationalDilationSemantics}\,X\);
  \item \(\code{externalization} :=
            \code{burnol\_blaschke\_detector\_externalization}\,
            S\,X\,\Obs.\code{externalIntro}\);
  \item \(\code{orthogonality} :=
            \code{burnol\_blaschke\_quotient\_orthogonality\_of\_intro}\,
            \Obs.\code{cokernelIntro}\);
  \item \(\code{admissibility} :=
            \code{burnol\_blaschke\_defect\_is\_admissible\_of\_intro}\,
            \Obs.\code{admissIntro}\).
\end{itemize}

The conclusion follows by
\Cref{thm:rh-from-inputs} applied to
\(\Obs.\code{bridgeRF01}\) and the assembled bundle.
\qedhere
\end{proof}

\subsection{The discharge of \(\code{bundledSubTargets\_of\_obstruction}\)}

The construction
\(\code{bundledSubTargets\_of\_obstruction}\) packages the five fields
of \code{BundledSubTargets} from a \(\Obs\) plus an explicit
detector \(S\) and a \code{DQObj} \(X\):
\begin{lstlisting}
noncomputable def bundledSubTargets_of_obstruction
    (Ω : Vol6FinalObstruction)
    (S : RKHSModelSpaceDetector burnolBlaschkeCondensedHilbertDefect)
    (X : DQObj) :
    Vol6.BundledSubTargets where
  detector := S
  rational :=
    Vol6.RationalDilationExternalization.canonicalRationalDilationSemantics X
  externalization :=
    Vol6.RationalDilationExternalization
      .burnol_blaschke_detector_externalization
      S X Ω.externalIntro
  orthogonality :=
    Vol6.QuotientOrthogonalityAdmissibility
      .burnol_blaschke_quotient_orthogonality_of_intro
      Ω.cokernelIntro
  admissibility :=
    Vol6.QuotientOrthogonalityAdmissibility
      .burnol_blaschke_defect_is_admissible_of_intro
      Ω.admissIntro
\end{lstlisting}
The \code{noncomputable} marker is forced by the use of
\code{burnol\_blaschke\_detector\_externalization}, which is
\code{noncomputable} on the Vol5 side because it operates on the
axiom type \code{NymanKernel}.  No proof obligation is incurred:
each field is a pure record-literal placement.

\subsection{Why no unconditional theorem is shipped}

\begin{theorem}[Honest non-delivery]
\label{thm:honest-non-delivery}
The synthesis module \code{Vol6.RHClassicalNewLanguage} \emph{does not}
ship a term

\begin{lstlisting}
theorem RH_classical_new_language : RH_classical
\end{lstlisting}

For Vol6 to ship such a term without a forbidden
\(\code{sorry}\)/\(\code{axiom}\)/\(\code{opaque}\) declaration, the
four fields of \(\code{Vol6FinalObstruction}\) would each need to be
inhabited from the Vol5 surface alone.  RF-01--RF-03 demonstrate this
is impossible: the Vol5 surface declares the carrier and the
admissibility atoms as
\(\code{axiom}\)/\(\code{opaque}\) without introduction rules.
\end{theorem}

\begin{proof}[Justification]
By inspection of the Vol5 surface
(\code{BurnolCore.lean}, \code{NoPhantomBlaschkeDefect.lean},
\code{NoPhantomConcreteAtoms.lean}), no term of any of
\(\code{burnolBlaschkeFactor.modelSpaceCarrier}\),
\(\code{BlaschkeDefectDualizable\,blaschkeDefectObject}\),
\(\code{BlaschkeDefectPolarized\,blaschkeDefectObject}\),
or \(\code{BlaschkeDefectRegularized\,blaschkeDefectObject}\)
is suppliable from the Vol5 closure (the carrier has no
constructor; the predicates have no introduction rule).  Vol6 cannot
modify Vol5 (forbidden) and cannot introduce new
\(\code{axiom}/\code{opaque}/\code{sorry}\) (forbidden).  Hence no
unconditional theorem is type-checkable, and the synthesis module
honestly omits it.
\qedhere
\end{proof}

\section{Success-criteria audit}
\label{sec:audit}

We now run, verbatim, the four success-criteria checks of
\code{next-steps.md}~\S``Success Criteria''.  Each is reported as
PASS, FAIL, or PARAMETERISED, with the precise Lean-level
artifact backing the verdict.

\subsection{Check 1: \texttt{lake build Vol6} passes}
\label{sec:check-1}

\begin{audit}[Lake build]
\label{audit:lake-build}
Run:
\begin{lstlisting}
cd lean && PATH="$HOME/.elan/bin:$PATH" lake build Vol6
\end{lstlisting}
Result: \textbf{PASS}.  The build produces

\begin{lstlisting}
Build completed successfully (2937 jobs).
\end{lstlisting}

with exit code 0, including the new module
\code{Vol6.RHClassicalNewLanguage}.
\end{audit}

\subsection{Check 2: no new \texttt{axiom}, \texttt{opaque}, \texttt{sorry}, \texttt{admit}}
\label{sec:check-2}

\begin{audit}[No new sorry/admit/axiom/opaque]
\label{audit:no-sorry}
Run:
\begin{lstlisting}
grep -RnE '\bsorry\b|\badmit\b|^axiom |^opaque ' lean/Vol6/
\end{lstlisting}
Result: \textbf{PASS}. All hits land in narrative \emph{comments} (e.g.\
``No \texttt{sorry}, \texttt{admit}, \dots\ introduced'' lines and
risk-flag descriptions); no proof-code occurrence is found.  The total
proof-code occurrence count is
\(\bigl(\code{sorry}, \code{admit}, \texttt{axiom}, \texttt{opaque}\bigr)
   = (0, 0, 0, 0)\).
\end{audit}

\subsection{Check 3: \texttt{\#print axioms} dependency list is clean}
\label{sec:check-3}

\begin{audit}[\(\code{\#print axioms}\) of the synthesis principal theorem]
\label{audit:print-axioms}
Run, against the principal theorem:
\begin{lstlisting}
#print axioms
  Vol6.RHClassicalNewLanguage.RH_classical_new_language_of_obstruction
\end{lstlisting}
Result: \textbf{PASS}. The output is
\begin{lstlisting}
'Vol6.RHClassicalNewLanguage.RH_classical_new_language_of_obstruction'
depends on axioms: [propext, Classical.choice, Quot.sound,
 Vol5.YonedaNymanTrackA.DQHom, Vol5.YonedaNymanTrackA.DQ_comp,
 Vol5.YonedaNymanTrackA.HardySubmodule, Vol5.YonedaNymanTrackA.NymanKernel,
 Vol5.YonedaNymanTrackA.blaschkeHardyImage,
 Vol5.YonedaNymanTrackA.burnolBlaschkeFactor]
\end{lstlisting}
The axiom set comprises (i) Lean/Mathlib foundational axioms
\(\code{propext},\,\code{Classical.choice},\,\code{Quot.sound}\)
and (ii) the pre-existing Vol5 axioms
\code{DQHom}, \code{DQ\_comp}, \code{HardySubmodule},
\code{NymanKernel}, \code{blaschkeHardyImage}, and
\code{burnolBlaschkeFactor}, all of which were already in the Vol5
inventory before Vol6 began. Crucially, the list contains
\emph{none} of: \code{RH\_classical}, \code{NymanDensity},
\code{ConjectureC\_reverse}, \code{InternalBlaschkeTriviality}, or
any Beurling-Nyman equivalence.
\end{audit}

\subsection{Check 4: model-space construction for OffCriticalZeroDefectKernel}
\label{sec:check-4}

\begin{audit}[Off-critical kernel via model-space construction]
\label{audit:off-critical}
Result: \textbf{PARAMETERISED-PASS}.
Paper~03 (\code{Vol6.OffCriticalDefectKernel}) ships
\(\code{off\_critical\_zero\_defect\_kernel\_of\_bridge}\) by
\emph{model-space construction}, not by contradiction from RH.
The proof route is:
\[
  s \;\xrightarrow{\;\alpha(s) = (s-1)/(s+1)\;}\; P_{s} = \{\alpha(s)\}
   \;\xrightarrow{\text{Paper 01}}\;
    \code{Nonempty}\,(\code{finiteBlaschkeModelSpace}\,P_{s}).\code{carrier}
   \;\xrightarrow{\text{bridge}}\;
   \code{Nonempty}\,\code{burnolBlaschkeFactor.modelSpaceCarrier}
\]
The first two arrows are unconditional (\Cref{sec:p01,sec:p03}).
The third is parameterised by \(\code{OffCriticalDefectKernelBridge}\).
The route never invokes RH or any RH-equivalent; the construction is
finite-Hardy-space, with the
canonical-lift dependency isolated in \(\Obs.\code{bridgeRF01}\).
\end{audit}

\subsection{Check 5: detector / externalization / admissibility for YonedaBlaschkeDetectorCalculus}
\label{sec:check-5}

\begin{audit}[Detector calculus by detector + externalization +
admissibility]
\label{audit:detector-calculus}
Result: \textbf{PARAMETERISED-PASS}.
Paper~06's
\(\code{yoneda\_blaschke\_detector\_calculus\_of\_subtargets}\) ships
the Target-2 calculus from the four-component bundle
\(\code{BundledSubTargets}\), whose components come precisely from
detector completeness (Paper~02), externalization (Paper~04),
quotient orthogonality (Paper~05A), and admissibility (Paper~05B). At
no point does the proof assume \(\code{NoPhantomDefect}\)
directly, the Beurling-Nyman criterion, or any RH-equivalent.
The four parameter dependencies are isolated in
\(\Obs.\code{externalIntro},\Obs.\code{cokernelIntro},\Obs.\code{admissIntro}\)
plus the Paper~02 explicit detector \(S\).
\end{audit}

\subsection{Check 6: no forbidden imports}
\label{sec:check-6}

Although not in the verbatim list of \code{next-steps.md}~\S
``Success Criteria'', the hard rules of the Vol6 contract include the
implicit constraint that no Vol6 module may import the
Beurling-Nyman criterion or Nyman density indirectly.  We audit
this here.

\begin{audit}[No forbidden Vol5 imports]
\label{audit:forbidden-imports}
Run:
\begin{lstlisting}
grep -RnE 'BurnolFactorization|YonedaNymanTrackA.Criterion'
  lean/Vol6/
\end{lstlisting}
Result: \textbf{PASS}.  The grep returns zero matches across all
twelve Vol6 modules.  No Vol6 module imports
\code{Vol5.YonedaNymanTrackA.BurnolFactorization} (which contains
the Nyman-density bridge \code{density\_iff\_blaschke\_trivial}) or
\code{Vol5.YonedaNymanTrackA.Criterion} (which contains the
Beurling-Nyman criterion).
\end{audit}

\subsection{Check 7: \(\code{Vol6FinalObstruction}\) is non-trivial}
\label{sec:check-7}

A formally vacuous obstruction (e.g.\ \code{True}) would also
satisfy all the previous checks, so we audit that
\(\Obs\) is \emph{actually} carrying mathematical content
proportional to RF-01--RF-04.

\begin{audit}[\(\Obs\) is non-trivial]
\label{audit:omega-nontrivial}
Result: \textbf{PASS}.
\(\code{Vol6FinalObstruction}\) has four fields:
(i) \code{bridgeRF01 : OffCriticalDefectKernelBridge}
    --- a \(\Type 1\)-valued function-typed structure;
(ii) \code{externalIntro : ExternalizationIntroductionRule \dots}
    --- a \(\Type 1\)-valued one-field structure;
(iii) \code{cokernelIntro : CanonicalQuotientOrthogonalityIntroduction}
    --- a \(\Prop\)-valued single-clause structure
    \emph{constructively equivalent} to
    \code{QuotientOrthogonalityInvisibility blaschkeDefectObject}
    by the audit theorem
    \code{quotient\_orthogonality\_invisibility\_iff\_introduction}; and
(iv) \code{admissIntro : CanonicalAdmissibilityIntroduction}
    --- a \(\Prop\)-valued three-field structure
    \emph{constructively equivalent} to
    \code{BlaschkeDefectAdmissible blaschkeDefectObject} by
    \code{admissibility\_intro\_iff}.

The two equivalence theorems prove that \(\Obs\) is exactly as
strong as the missing Vol5 surface, not stronger.  The structure is
therefore non-vacuous; in particular it is not provable in vol6
without inhabiting at least one of RF-01--RF-03.
\end{audit}

\subsection{Combined verdict on the five checks}

\begin{table}[h]
\centering
\begin{tabular}{@{}l l l@{}}
\toprule
Check & Result & Backing artifact \\
\midrule
1. \code{lake build Vol6}        & PASS               &
  \Cref{audit:lake-build} \\
2. No new sorry/admit/axiom/opaque & PASS             &
  \Cref{audit:no-sorry} \\
3. \(\code{\#print axioms}\) clean & PASS             &
  \Cref{audit:print-axioms} \\
4. OffCriticalZeroDefectKernel by model space & PARAMETERISED-PASS &
  \Cref{audit:off-critical} \\
5. YonedaBlaschkeDetectorCalculus by detector+ext+admiss & PARAMETERISED-PASS &
  \Cref{audit:detector-calculus} \\
\bottomrule
\end{tabular}
\caption{Success-criteria audit summary.}
\label{tab:success-criteria}
\end{table}

\section{Honest verdict}
\label{sec:verdict}

\paragraph{Closure status.}
Classical RH \emph{does not} close unconditionally inside Vol6.
The synthesis principal theorem is parameterised over
\(\code{Vol6FinalObstruction}\) plus an
\(\code{RKHSModelSpaceDetector}\) for
\code{burnolBlaschkeCondensedHilbertDefect} and a \code{DQObj}.

\paragraph{Single residual barrier.}
\(\code{Vol6FinalObstruction}\) is the unique structure whose
inhabitation closes the Vol7 work.  Its four fields are
each transparent records (no opaque/axiom dependency beyond the Vol5
inheritance), but inhabitation requires Vol7 to expose new Vol5-side
introduction rules, all four of which are mathematically standard
constructions blocked only by the absence of constructors at the
Vol5 surface.

\paragraph{What we did not assume.}
The synthesis module \code{Vol6.RHClassicalNewLanguage} does not
import \code{Vol5.YonedaNymanTrackA.BurnolFactorization}, does not
import \code{Vol5.YonedaNymanTrackA.Criterion} (the
Beurling-Nyman-equivalent criterion), does not introduce a
\(\code{NymanDensity}\) hypothesis, and does not treat
\(\code{InternalBlaschkeTriviality}\) as an assumption.  The
\(\code{\#print axioms}\) check confirms this.

\paragraph{Why this is the honest deliverable.}
Vol6's contract is to assemble what the
\code{NoPhantomLanguageBreakthrough} payload \emph{needs} from
honest finite-rank constructions.  The finite-rank content (Papers
01--02) is closed.  The canonical-lift content (Papers 03--06) is
parameterised over the precise Vol5 introduction rules that would,
if available, close the entire chain.  Anything more than this
without modifying Vol5 would require introducing forbidden
\(\code{axiom}\)/\(\code{opaque}\) declarations, which would amount
to \emph{pretending} we have closed the obstruction.  The
\(\Obs\) structure is therefore the maximally-honest
deliverable.

\section{Vol7 roadmap}
\label{sec:vol7-roadmap}

Vol7's job is to inhabit \(\code{Vol6FinalObstruction}\) and apply
\Cref{thm:synthesis-principal}.  We sketch the four bridge
constructions Vol7 must produce.

\subsection{Vol7 bridge \#1 --- the canonical-carrier transport}

\paragraph{Target.}
\(\code{OffCriticalDefectKernelBridge}\), i.e.\ a function
\[
  \forall (P : \code{FiniteBlaschkePacket}),\;
   \code{Nonempty}\,P.\code{ZeroIndex}
   \to (\code{finiteBlaschkeModelSpace}\,P).\code{carrier}
   \to \code{burnolBlaschkeFactor.modelSpaceCarrier}.
\]

\paragraph{Strategy.}
Either (a) replace the \code{axiom burnolBlaschkeFactor} with a
concrete vol7 definition that ties the carrier to a quotient of
\(\code{HardySpaceH2}\) by \(\code{blaschkeHardyImage}\), or (b)
extend the Vol5 surface (in vol7 only, not vol5) with a
\(\code{constructor}\) lemma that builds the canonical carrier from
finite-packet model spaces.  Strategy (a) is preferred because it
makes the carrier concrete; the cost is a non-trivial vol7
\code{HardyModelSpace} development.

\subsection{Vol7 bridge \#2 --- the externalization introduction rule}

\paragraph{Target.}
\(\code{ExternalizationIntroductionRule}\,
   \code{burnolBlaschkeCondensedHilbertDefect}\), i.e.\ for every
\(\DQ\)-presheaf \(F\) that is a \(\code{RationalDilationRepresentable}\),
the predicate
\(\code{CondensedHilbertDefectDetectedByRationalDilation}\) holds.

\paragraph{Strategy.}
The natural route is via \code{Vol5.YonedaNymanTrackA.NoPhantomConcreteAtoms},
which proves the rigidity \(\code{DQObj} = \code{NymanKernel}\)
definitionally.  Vol7 should expose a vol5-side characterization
that the opaque \code{DefectDetectedByRationalDilationRepresentable}
predicate is provably equivalent to a transparent statement of
``existence of a detector functional with nonzero pairing''.  This
is RF-02's resolution.

\subsection{Vol7 bridge \#3 --- the cokernel introduction rule}

\paragraph{Target.}
\(\code{CanonicalQuotientOrthogonalityIntroduction}\), i.e.\ a proof
that \(\code{YonedaNymanRepresentableImage}\,F\) implies
\(\neg\,\code{DefectDetectedByRationalDilationRepresentable}\,
   \code{blaschkeDefectObject}\,F\).

\paragraph{Strategy.}
Combine the previous bridge with the no-phantom rigidity theorem
\(\code{no\_phantom\_blaschke\_defect\_rigidity}\) of
\code{NoPhantomConcreteAtoms.lean}.  Both directions are already
proved on the Vol5 side modulo the opaque predicate; lifting
through the bridge of \S\ref{sec:vol7-roadmap}.2 closes this.

\subsection{Vol7 bridge \#4 --- the three admissibility atoms}

\paragraph{Target.}
\(\code{CanonicalAdmissibilityIntroduction}\), i.e.\ proofs of
\(\code{BlaschkeDefectDualizable},
  \code{BlaschkeDefectPolarized},
  \code{BlaschkeDefectRegularized}\) for
\(\code{blaschkeDefectObject}\).

\paragraph{Strategy.}
At the finite-rank level all three atoms are theorems of standard
finite-dimensional Hilbert-module theory:
\begin{itemize}
  \item Dualisability: a finite-dimensional Hilbert space is a
        dualisable module, with snake identities provable by trace
        formulas.
  \item Polarisation: the Hardy inner product restricted to the
        finite-rank model space \(\KB^{(P)}\) is a Hermitian form
        admitting the polarisation identity.
  \item Regularised convergence: the resolvent family
        \((R(z))_{z \in \rho(T_{B})}\) of the truncated
        multiplication operator on \(\KB^{(P)}\) converges in the
        regularised topology by the spectral theorem in finite
        dimension.
\end{itemize}
Lifting from the finite-rank level to the canonical condensed
defect requires a Vol7 completion/colimit argument; the precise
Vol5 surface to extend is the
\code{CondensedHilbertDefect} layer in
\code{Vol5.YonedaNymanTrackA.CondensedHilbertDefect}.

\subsection{What Vol7 ships}

Vol7 ships:

\begin{itemize}[leftmargin=2em]
  \item Four Lean theorems closing each bridge above
        (\code{vol7\_bridge\_RF01}, \code{vol7\_bridge\_external},
        \code{vol7\_bridge\_cokernel},
        \code{vol7\_bridge\_admiss}).
  \item Their assembly
        \code{vol7\_finalObstruction : Vol6FinalObstruction}.
  \item The unconditional theorem
        \(\code{RH\_classical\_new\_language : RH\_classical} :=
            \code{RH\_classical\_new\_language\_of\_obstruction}\,
              \code{vol7\_finalObstruction}\,S\,X\)
        with explicit witnesses \(S\) and \(X\) constructed in
        Vol7.
\end{itemize}

\subsection{Estimated cost of each Vol7 bridge}

We sketch a rough order-of-magnitude estimate of the formalisation cost
of each bridge.  All estimates are in module-LOC (lines of Lean code),
loosely benchmarked against the corresponding Vol6 papers.

\begin{table}[h]
\centering
\begin{tabular}{@{}l r r l@{}}
\toprule
Bridge & Estimated LOC & Mathlib dependency & Mathematical content \\
\midrule
\#1 (RF-01)        &  600 & Mathlib.Hardy        & Concrete \(\HtH\)
                                                  with model-space
                                                  quotient \\
\#2 (Externalization) & 400 & Mathlib.CategoryTheory.Yoneda
                                                  & Detector $\to$ Yoneda
                                                  rep equivalence \\
\#3 (Cokernel)        & 300 & Mathlib.HilbertOrthogonal & Orthogonal
                                                  complement of
                                                  \(\code{blaschkeHardyImage}\) \\
\#4 (Admissibility)   & 800 & Mathlib.Resolvent     & Spectral resolvent
                                                  + finite-rank
                                                  trace formulas \\
\midrule
Total                 & 2100 &                       & \\
\bottomrule
\end{tabular}
\caption{Estimated Vol7 module sizes by bridge.}
\label{tab:vol7-budget}
\end{table}

The estimates assume the existing Vol6 surface remains read-only and
that Vol7 is layered on top in the same Lake-package style as Vol6.
Bridge \#1 (RF-01) is the largest single piece because it requires
porting Mathlib's Hardy space machinery and producing a concrete
\(\code{burnolBlaschkeFactor}\) instance.  Bridge \#4
(admissibility) is the second largest because it requires the
full spectral-theorem machinery for the truncated multiplication
operator.

\subsection{Comparison with the Beurling-Nyman route}

The reader familiar with the classical Beurling-Nyman programme may
ask: why not just close the Vol7 obstruction by appealing to
\(\code{NymanDensity}\)?  The answer is in two parts.  First, the
Vol6 hard rules forbid this: \(\code{NymanDensity}\) is a
proposition equivalent to RH, and assuming it would amount to
assuming RH.  Second, the entire point of the No-Phantom Language
is that the contradiction is moved into HoTT/Yoneda visibility:
\emph{no admissible Yoneda-invisible defect can exist}.  This is a
strictly stronger meta-theoretic stance than ``Nyman-density implies
RH'', because it produces RH from a structural rigidity of the
condensed Hilbert framework rather than from an analytic density
estimate.  Vol7's job is to close the obstructions
\emph{constructively}, not to replace them with an RH-equivalent
density.

\subsection{Risk assessment for Vol7}

We score each bridge on two axes: \emph{mathematical risk} (is the
underlying mathematics standard?) and \emph{formalisation risk} (is
the Lean translation routine?).  See \Cref{tab:vol7-risk}.

\begin{table}[h]
\centering
\begin{tabular}{@{}l c c l@{}}
\toprule
Bridge & Math risk & Lean risk & Comment \\
\midrule
\#1 (RF-01)            & low  & medium &
   Hardy theory is classical; Mathlib coverage of
   \(\HtH\) is partial. \\
\#2 (Externalization)   & low  & low &
   Yoneda lemma is mathlib-native. \\
\#3 (Cokernel)          & low  & low &
   Orthogonal-complement semantics are mathlib-native. \\
\#4 (Admissibility)     & medium & high &
   Regularised resolvent convergence; care with non-norm topologies. \\
\bottomrule
\end{tabular}
\caption{Vol7 risk assessment.}
\label{tab:vol7-risk}
\end{table}

\subsection{Vol7 audit will mirror Vol6}

Vol7 will inherit the same hard-rule contract as Vol6: no
\(\code{sorry}\), no \(\code{admit}\), no new
\(\code{axiom}\) outside the inherited Vol5/Mathlib surface, no new
\(\code{opaque}\), no Nyman density or Beurling-Nyman, and no
RH-equivalent.  The Vol7 \(\code{\#print axioms}\) audit is
expected to produce the \emph{same} output as
\Cref{audit:print-axioms}: only the Lean foundational axioms and
the inherited Vol5 surface.  Concretely Vol7 should \emph{not}
introduce new axiom types and should aim to retire \emph{at least}
RF-01's \code{burnolBlaschkeFactor} axiom by replacing it with a
concrete definition.

\section{Comparison with the prior Volumes}
\label{sec:vol-comparison}

\subsection{Vol4 vs.\ Vol5 vs.\ Vol6}

The HoTT/Yoneda Riemann hypothesis programme has progressed through
four volumes (Vol4 forward).  Each volume's contribution can be
summarised by what it shipped at its synthesis level:

\begin{itemize}[leftmargin=2em]
  \item \textbf{Vol4 (Track-A formulation).}  Introduced the
        formal target \code{RH\_classical} as an HoTT-internal
        proposition, established the V5a Hilbert/Pólya
        compatibility ladder, and identified the conditional shape
        of an RH proof through the spectral-side reduction.  Vol4
        ships axioms \code{Infra1\_HP\_collapse},
        \code{Infra2\_zeta\_compatibility},
        \code{Infra4\_spectral}, etc., which Vol5 inherits.
  \item \textbf{Vol5 (No-Phantom Language).}  Reformulated RH as
        the conditional theorem
        \code{RH\_classical\_of\_no\_phantom\_language\_breakthrough},
        introducing the Yoneda/Blaschke detector calculus +
        off-critical defect kernel as the canonical payload.
        Vol5 closes the conditional theorem unconditionally; the
        payload is open.  Vol5 \emph{does not} construct the
        canonical \(\code{burnolBlaschkeFactor}\); it ships it as
        an axiom.
  \item \textbf{Vol6 (Honest Reduction).}  The present sprint:
        produces the finite-rank machinery (Papers 01--02), bridges
        it to the Vol5 surface (Papers 03, 06), and reduces the
        residual obligation to a single \(\Type 1\)-valued
        obstruction \(\code{Vol6FinalObstruction}\).  Vol6 ships no
        new axiom; it parameterises over the four named
        introduction rules.
\end{itemize}

\subsection{What changed conceptually}

Conceptually, the move from Vol4 to Vol6 is the move from a
\emph{spectral} formulation of RH (Hilbert--Pólya, ECond, the V5a
ladder) to a \emph{representational} formulation (the canonical
Burnol/Blaschke defect, condensed Hilbert detection, rational-dilation
Yoneda invisibility).  This move is not a strengthening of the proof
but a reformulation: it makes the obstruction visible at the level of
type-theoretic data (introduction rules) rather than at the level of
analytic estimates (zero-counting).

\subsection{What is the same}

What is the same is the underlying axiom inventory.  Vol6 inherits
all of Vol5's axioms and adds none of its own; the
\(\code{\#print axioms}\) audit
(\Cref{audit:print-axioms}) confirms this.  In particular Vol6 does
not strengthen \code{RH\_classical}, it only re-shapes the route
through the No-Phantom Language and isolates the precise residual
obstruction.

\section{Detailed walk-through of \(\code{RH\_classical\_new\_language\_of\_obstruction}\)}
\label{sec:walkthrough}

For a reader who wishes to verify the Lean proof line-by-line, we
walk through the synthesis principal theorem in detail.

\subsection{Preconditions}

The hypothesis triple is \((\Obs, S, X)\):
\begin{itemize}[leftmargin=2em]
  \item \(\Obs : \code{Vol6FinalObstruction}\) provides the
        four named introduction-rule fields.
  \item \(S : \code{RKHSModelSpaceDetector}\,
            \code{burnolBlaschkeCondensedHilbertDefect}\)
        is an \emph{arbitrary} (non-vacuous or vacuous) detector for
        the canonical condensed Hilbert defect.  It is data, not
        opaque-Prop discharge.
  \item \(X : \code{DQObj}\) is an \emph{arbitrary} \(\DQ\)-object.
        \code{DQObj} is itself an \code{axiom Type}
        in Vol5, but Vol5 also ships
        \code{NymanKernel : Type} (an axiom),
        and \code{DQObj := NymanKernel} definitionally.
        Producing an inhabitant \(X\) is therefore at the same
        level of difficulty as inhabiting \code{NymanKernel},
        a separate Vol7 task.
\end{itemize}

\subsection{Step 1: build the off-critical kernel}

From \(\Obs.\code{bridgeRF01} : \code{OffCriticalDefectKernelBridge}\)
we apply Paper~03's
\code{off\_critical\_zero\_defect\_kernel\_of\_bridge} to obtain
\(\code{OffCriticalZeroDefectKernel}\).  This step uses
Paper~01's \code{finite\_packet\_model\_space\_nonempty} via the
singleton-packet construction.  No opaque-Prop discharge is needed.

\subsection{Step 2: build the externalization}

From \(\Obs.\code{externalIntro}\) and the explicit
\((S, X)\) we apply Paper~04's
\code{burnol\_blaschke\_detector\_externalization} to obtain
\(\code{RKHSDetectorExternalizationToRationalRepresentables}\,K\,S\,
   (\code{canonicalRationalDilationSemantics}\,X)\)
where \(K = \code{burnolBlaschkeCondensedHilbertDefect}\).

\subsection{Step 3: build the orthogonality}

From \(\Obs.\code{cokernelIntro}\) we apply Paper~05A's
\code{burnol\_blaschke\_quotient\_orthogonality\_of\_intro} to obtain
\(\code{QuotientOrthogonalityInvisibility}\,
   \code{blaschkeDefectObject}\).  No detector or rational data is
needed.

\subsection{Step 4: build the admissibility}

From \(\Obs.\code{admissIntro}\) we apply Paper~05B's
\code{burnol\_blaschke\_defect\_is\_admissible\_of\_intro} to obtain
\(\code{BurnolBlaschkeDefectIsAdmissible}\), which is an abbreviation
for \(\code{AdmissibleCondensedHilbertDefect}\,
   \code{burnolBlaschkeCondensedHilbertDefect}\).

\subsection{Step 5: assemble the bundle}

We package
\((S, \code{canonicalRationalDilationSemantics}\,X,
   \text{externalization},
   \text{orthogonality}, \text{admissibility})\)
into a \(\code{Vol6.BundledSubTargets}\) record.  This is exactly the
\(\code{bundledSubTargets\_of\_obstruction}\) construction.

\subsection{Step 6: assemble the calculus}

Paper~06's
\(\code{assembleYonedaBlaschkeDetectorCalculus}\) consumes
the bundle and produces
\(\code{YonedaBlaschkeDetectorCalculus}\) by a single record
literal.

\subsection{Step 7: assemble the breakthrough payload}

We pair the calculus with the off-critical kernel of Step~1 to
form the
\(\code{NoPhantomLanguageBreakthrough}\) record.

\subsection{Step 8: apply the Vol5 master theorem}

Vol5's
\(\code{RH\_classical\_of\_no\_phantom\_language\_breakthrough}\)
takes the breakthrough payload and emits \(\RH\).  Done.

\subsection{Audit summary of the eight steps}

\begin{table}[h]
\centering
\small
\begin{tabular}{@{}r l l@{}}
\toprule
Step & Lemma applied & Source paper \\
\midrule
1 & \code{off\_critical\_zero\_defect\_kernel\_of\_bridge}            & P03 \\
2 & \code{burnol\_blaschke\_detector\_externalization}                 & P04 \\
3 & \code{burnol\_blaschke\_quotient\_orthogonality\_of\_intro}         & P05A \\
4 & \code{burnol\_blaschke\_defect\_is\_admissible\_of\_intro}          & P05B \\
5 & \code{bundledSubTargets\_of\_obstruction} (record literal)          & P07 \\
6 & \code{assembleYonedaBlaschkeDetectorCalculus} (record literal)     & P06 \\
7 & \code{NoPhantomLanguageBreakthrough} (record literal)              & V5 \\
8 & \code{RH\_classical\_of\_no\_phantom\_language\_breakthrough}      & V5 \\
\bottomrule
\end{tabular}
\caption{The eight-step composition.}
\label{tab:eight-step}
\end{table}

\subsection{Audit theorem: the eight-step composition is total}

There is no `partiality' in the eight-step composition: every step is
a definitional record-literal construction or an application of a
lemma that is proved unconditionally up to the obstruction inputs.
The composition fails if and only if at least one of the four
inhabiting hypotheses
\(\Obs.\code{bridgeRF01},\; \Obs.\code{externalIntro},\;
   \Obs.\code{cokernelIntro},\; \Obs.\code{admissIntro}\) is missing.

\subsection{Independence of the four obstructions}

We have established that the four fields of \(\Obs\) are
\emph{individually} required.  Are they \emph{jointly} required?
Yes: each addresses a distinct opaque/axiom Vol5 atom, and no
field's content is derivable from the other three on the Vol5
surface.

\begin{itemize}[leftmargin=2em]
  \item \code{bridgeRF01} addresses
        \(\code{burnolBlaschkeFactor.modelSpaceCarrier}\).
  \item \code{externalIntro} addresses
        \(\code{DefectDetectedByRationalDilationRepresentable}\).
  \item \code{cokernelIntro} addresses the same opaque predicate
        from a different direction (the \emph{negation} for
        Yoneda-image inputs).
  \item \code{admissIntro} addresses the three opaque admissibility
        atoms.
\end{itemize}

The third and second appear redundant at first glance (both deal
with \(\code{DefectDetectedByRationalDilationRepresentable}\)), but
they are not: one is the \emph{introduction} (positive) form for
representable-detection and the other is the \emph{negation}
(refutation) form for Yoneda-image inputs.  The opaque predicate
admits no derived eliminator from which the negation could be
inferred from the positive introduction; thus both are needed
independently.

\section{Frequently anticipated objections}
\label{sec:objections}

\paragraph{Objection 1: ``Isn't \(\Obs\) just RH in disguise?''}
No.  The fields of \(\Obs\) are each \emph{constructively}
weaker than RH: each is a transparent introduction rule for an
opaque Vol5 atom, and the inhabitation of these rules requires
finite-rank/spectral machinery rather than zero-counting on the
critical line.  More concretely, \(\Obs\) does not entail
\(\code{NymanDensity}\) (which would be RH-equivalent), nor
the Beurling criterion, nor any density of a Mellin image.

\paragraph{Objection 2: ``Why not weaken the Vol5 surface?''}
The Vol6 hard rules forbid modifying the Vol5 surface.  That is
the right contract for synthesis: Vol7 is allowed to extend the
Vol5 surface (or replace its axioms with concrete definitions)
\emph{provided it does not introduce new axioms or RH-equivalents}.

\paragraph{Objection 3: ``Why not just inhabit the obstruction with
\(\code{Classical.choice}\)?''}
Because that would amount to introducing a new axiom: Lean's
\code{Classical.choice} is already in the inherited axiom inventory,
but using it to inhabit the four obstruction fields would require
a witness in each case, and the obstruction \(\Obs\)'s
fields are \emph{universally quantified function types} or \(\Prop\)s
that are not provable from \code{Classical.choice} alone.  In
particular \code{bridgeRF01} requires a function out of a non-empty
\(\Type\) into a non-empty \(\Type\); \code{Classical.choice} gives no
such function unless we already know both types are inhabited and
can be related.

\paragraph{Objection 4: ``Could Vol6 have closed unconditionally
if Paper~05 used Mathlib's \(\HtH\)?''}
No.  Mathlib's Hardy space is defined on the unit circle as a
subspace of \(L^{2}(\partial\D)\); the Vol5
\code{burnolBlaschkeFactor.modelSpaceCarrier} is an
\code{axiom Type} unrelated to Mathlib's construction.  The
canonical bridge between the two is exactly the missing RF-01
content; Vol7 must supply it.

\paragraph{Objection 5: ``What if a future RH proof uses different
language?''}
Vol7 may certainly take a different route.  But Vol6's value is
that it isolates the precise residual obstruction in the
No-Phantom Language framework, which any subsequent Vol6+
extension must address.  If Vol7 chooses an entirely different
route (e.g.\ a return to Beurling-Nyman with a constructive
density proof), then \(\Obs\) becomes irrelevant, but the
Vol5 conditional theorem and Vol6's
\code{RH\_classical\_new\_language\_of\_payload} wrapper remain
valid as alternative entry points.

\section{Conclusion}
\label{sec:conclusion}

Vol6 reduces classical RH to a single \(\Type 1\)-valued obstruction
\(\code{Vol6FinalObstruction}\), with four named fields whose
inhabitation is the precise Vol5-surface extension required to close
the No-Phantom Language route.  All five \code{next-steps.md}
success-criteria checks pass, with the unconditional-closure check
honestly re-stated as the \(\Obs\) obstruction.  No new
\(\code{sorry}\), \(\code{admit}\), \(\code{axiom}\), or
\(\code{opaque}\) is introduced; the \(\code{\#print axioms}\)
output is clean of RH, Nyman density, Beurling-Nyman, and any
RH-equivalent.  The canonical Vol7 deliverable is the four bridge
theorems sketched in \Cref{sec:vol7-roadmap}.

\appendix

\section{Reproduced Lean signatures}
\label{appx:lean-sigs}

For the reader's convenience we reproduce the eight key Lean
signatures that drive the synthesis.  All declarations are
\(\code{sorry}/\code{admit}/\code{axiom}/\code{opaque}\)-free.

\subsection{The Vol5 master theorem (reused, not modified)}

\begin{lstlisting}
-- Vol5.YonedaNymanTrackA.NoPhantomLanguage, line 137
theorem RH_classical_of_no_phantom_language_breakthrough
    (h : NoPhantomLanguageBreakthrough) :
    RH_classical
\end{lstlisting}

\subsection{The Paper 03 conditional theorem}

\begin{lstlisting}
-- Vol6.OffCriticalDefectKernel
theorem off_critical_zero_defect_kernel_of_bridge
    (bridge : OffCriticalDefectKernelBridge) :
    OffCriticalZeroDefectKernel
\end{lstlisting}

\subsection{The Paper 06 conditional assembly}

\begin{lstlisting}
-- Vol6.YonedaBlaschkeDetectorCalculus
def yoneda_blaschke_detector_calculus_of_subtargets
    (B : BundledSubTargets) :
    YonedaBlaschkeDetectorCalculus
\end{lstlisting}

\subsection{The Vol6 final obstruction}

\begin{lstlisting}
-- Vol6.RHClassicalNewLanguage
structure Vol6FinalObstruction : Type 1 where
  bridgeRF01    : OffCriticalDefectKernelBridge
  externalIntro : ExternalizationIntroductionRule
                    burnolBlaschkeCondensedHilbertDefect
  cokernelIntro : CanonicalQuotientOrthogonalityIntroduction
  admissIntro   : CanonicalAdmissibilityIntroduction
\end{lstlisting}

\subsection{The synthesis wrapper}

\begin{lstlisting}
-- Vol6.RHClassicalNewLanguage
theorem RH_classical_new_language_of_payload
    (payload : NoPhantomLanguageBreakthrough) :
    RH_classical
\end{lstlisting}

\subsection{The substantive deliverable}

\begin{lstlisting}
-- Vol6.RHClassicalNewLanguage
theorem RH_classical_new_language_of_inputs
    (P03 : OffCriticalDefectKernelBridge)
    (P06 : Vol6.BundledSubTargets) :
    RH_classical
\end{lstlisting}

\subsection{The bundled-sub-targets-of-obstruction record}

\begin{lstlisting}
-- Vol6.RHClassicalNewLanguage
noncomputable def bundledSubTargets_of_obstruction
    (Ω : Vol6FinalObstruction)
    (S : RKHSModelSpaceDetector burnolBlaschkeCondensedHilbertDefect)
    (X : DQObj) :
    Vol6.BundledSubTargets
\end{lstlisting}

\subsection{The principal synthesis theorem}

\begin{lstlisting}
-- Vol6.RHClassicalNewLanguage
theorem RH_classical_new_language_of_obstruction
    (Ω : Vol6FinalObstruction)
    (S : RKHSModelSpaceDetector burnolBlaschkeCondensedHilbertDefect)
    (X : DQObj) :
    RH_classical
\end{lstlisting}

\section{Verbatim \texttt{\#print axioms} output}
\label{appx:print-axioms}

\begin{lstlisting}
'Vol6.RHClassicalNewLanguage.RH_classical_new_language_of_obstruction'
depends on axioms: [propext,
 Classical.choice,
 Quot.sound,
 Vol5.YonedaNymanTrackA.DQHom,
 Vol5.YonedaNymanTrackA.DQ_comp,
 Vol5.YonedaNymanTrackA.HardySubmodule,
 Vol5.YonedaNymanTrackA.NymanKernel,
 Vol5.YonedaNymanTrackA.blaschkeHardyImage,
 Vol5.YonedaNymanTrackA.burnolBlaschkeFactor]
\end{lstlisting}

\end{document}
